
\subsection{Overview of System Approach}

The suggested approach uses a closed-loop antenna alignment mechanism to dynamically orient the directional antenna based on RSSI feedback received during the process. In essence, the objective is to increase the signal strength at the receiving node by continuously adjusting the orientation of the antenna depending on the current wireless channel state.

From the point of view of functionality, the system works as a simple feedback control loop comprised of three main components: signal sensing, decision making, and actuator control. RSSI values are continuously measured at the receiving node, then processed by a specific alignment algorithm, and converted into motor commands for rotating the antenna. Every change in the position of the antenna results in further sensing of RSSI signals, creating a full cycle in a control loop.

Unlike open-loop antenna alignment or manual orientation of an antenna, the suggested solution does not involve any initial knowledge about the direction of the transmitter or particular models of the channel. Instead, all alignment actions are made based on feedback alone, thus allowing the use of propagation channels whose characteristics are unknown or constantly changing.

To estimate the effectiveness of adaptive antenna alignment, several different orientation algorithms can be used, including conventional approaches and reinforcement learning (RL).

Baseline strategies include:
\begin{itemize}
    \item Angular scanning with maximum RSSI, where the antenna is rotated in order to determine the optimal angle with respect to the maximum RSSI.
    \item Hill climbing based on RSSI, which consists of small adjustments made to the antenna according to RSSI gradients.
\end{itemize}

Such approaches form the baseline performance measures that differ in their precision of alignment, speed, resource consumption, and robustness against noise, since exhaustive scanning is expected to yield optimal orientation in the static case, whereas hill climbing will suffer due to the possibility of local optimum.

The reinforcement learning approach represents a method of making decisions for achieving certain goals by interacting with the environment. Instead of performing many scans, the RL system will make the necessary adjustments based on reward received by comparing RSSI values. Thus, it becomes possible to compare a heuristic versus a learning-based control in similar conditions.

Several decisions regarding design were made with the goal of increasing feasibility, repeatability, and embeddability.

The ESP32 microcontroller board was selected due to its built-in WiFi radio, ability to read RSSI from hardware, sufficient computing power, and accessibility. Ability to handle wireless communication, signal reading, and control actions on a single and affordable piece of hardware is crucial to integration simplicity and reduction of dependencies.

A stepper motor is utilized for antenna movement due to the possibility to accurately control the angle in which the antenna is positioned relative to the target node. As opposed to servo motors, stepper motors allow achieving deterministic movement without any additional sensors, thus simplifying the mechanical part of the design and the control algorithm. Discrete rotational angle steps make it easy to align this feature with discrete action space needed for light reinforcement learning.

Reinforcement learning model is intentionally designed to work within the discrete state and action space. Such approach helps avoid excessive resource consumption required by continuous-state RL and deep learning algorithms. They simply aren't feasible when used on constrained hardware, whereas tabular learning methods are easier to implement, predictable in terms of memory requirements, and robust against noisy measurements.

It is important to note that RL in this design was treated as a way to optimize the control method, but not as a standalone optimizer. Control works without any training, and all baseline methods can be independently run and tested.